\section{Assembly Wrapping: Mechanism and Integration}\label{sec:wrapping}
\subsection{The Same Bytes Can Have Two Instruction Interpretations}
Assembly wrapping changes how executable bytes are arranged and interpreted.
It uses the variable length of x86 and x64 instructions: decoding the same
bytes from different starting positions can produce different instruction
sequences. In the existing wrapping technique, bytes belonging to the
intended computation are placed within the immediate-operand field of a
larger instruction~\cite{assembly_wrapping}. An \emph{immediate operand} is
a constant stored directly in an instruction's encoding. A decoder starting
at the outer instruction reads that field as a constant; execution entering
the embedded sequence reads it as instructions instead.

Figure~\ref{fig:wrapping_views} illustrates the distinction without relying
on a particular encoding. The \emph{container} is the outer instruction
whose operand holds the embedded bytes. The execution path connects the
embedded sequences in their intended order, whereas a static listing
following the container boundaries shows a different instruction sequence.
The bytes need not be decrypted to obtain that second interpretation.
This is a representation technique based on overlapping instructions, not
a new form of encryption or a guarantee that an analyzer cannot recover
the executed code.

\begin{figure}[t]
\centering
\begin{tikzpicture}[font=\small]
\node[anchor=east] at (0,0.5) {Stored bytes};
\draw[fill=gray!15] (0.15,0.15) rectangle (2.7,0.85);
\draw[fill=blue!8] (2.7,0.15) rectangle (7.4,0.85);
\node at (1.425,0.5) {Outer opcode};
\node at (5.05,0.5) {Immediate-operand field};
\node[anchor=east] at (0,1.65) {Outer view};
\node[draw,align=center,text width=6.9cm] at (3.775,1.65)
  {One instruction containing a constant operand};
\node[anchor=east] at (0,-0.9) {Execution view};
\node[draw,fill=blue!8,align=center,text width=4.1cm] (inside)
  at (5.05,-0.9) {Embedded instruction sequence};
\draw[-{Latex[length=2mm]},thick] (2.7,0.15) -- (2.7,-0.9) -- (inside.west);
\end{tikzpicture}
\caption{Conceptual views of a wrapped region. Starting at the outer
opcode interprets the operand as a constant; entering the embedded
sequence interprets some of those same bytes as instructions.
The drawing is not to scale and omits inter-region control transfers.}
\label{fig:wrapping_views}
\end{figure}

\subsection{What Zone Adds and What the Transformation Requires}
The contribution of this work is the integration and assessment of an
existing technique, with implementations for both architectures. Zone
starts from a prepared instruction stream, groups supported instructions
into wrapped regions, preserves the required execution relationships, and
updates the supported location-dependent references after layout. Some
instructions or local layouts cannot be placed in the main container
representation; the implementation retains these in separate, unwrapped
regions. An output may therefore contain both wrapped and ordinary code.

The x86 and x64 implementations share this purpose but use different
encodings and addressing rules. In particular, the available immediate
operand has a different width, affecting which instruction groups fit.
The x86 implementation was introduced after the first benign benchmark
snapshot and is evaluated under its own recorded version.

Data preparation remains necessary. Wrapping does not determine whether
an apparent instruction is actually an embedded string or table. If that
earlier interpretation is wrong, construction may be refused or the output
may fail when executed. Successful relocation checks likewise cover the
references the engine supports, not every possible address computation.

\subsection{What the Evaluation Can Establish}
Three questions follow from this mechanism: can Zone construct an output
for the input, does that output pass its behavioral test, and what do static
scanners report for the produced bytes? These questions have different
denominators. A scan cannot evaluate an unbuilt output, and an absence of
malicious verdicts does not demonstrate that a constructed output works.
The following experiments therefore report build availability, functional
outcomes, and scan results separately, including the failures on the full
102-payload corpus rather than only the successful benign examples.
